\documentclass[10pt,a4paper]{amsart}
\usepackage{amsmath,amssymb,amsthm,graphicx}
\usepackage{verbatim}
\newenvironment{reference}{\comment}{\endcomment}
%\newenvironment{reference}{}{}
\newenvironment{slogan}{\comment}{\endcomment}
\newenvironment{history}{\comment}{\endcomment}
% For commutative diagrams we use Xy-pic
\usepackage[all]{xy}

% We use 2cell for 2-commutative diagrams.
\xyoption{2cell}
% Theorem environments.
%
\theoremstyle{plain}
\newtheorem{theorem}[subsection]{Theorem}
\newtheorem{proposition}[subsection]{Proposition}
\newtheorem{lemma}[subsection]{Lemma}

\theoremstyle{definition}
\newtheorem{definition}[subsection]{Definition}
\newtheorem{example}[subsection]{Example}
\newtheorem{exercise}[subsection]{Exercise}
\newtheorem{situation}[subsection]{Situation}

\theoremstyle{remark}
\newtheorem{remark}[subsection]{Remark}
\newtheorem{remarks}[subsection]{Remarks}

\numberwithin{equation}{subsection}

% Macros
%
\def\lim{\mathop{\mathrm{lim}}\nolimits}
\def\colim{\mathop{\mathrm{colim}}\nolimits}
\def\Spec{\mathop{\mathrm{Spec}}}
\def\Hom{\mathop{\mathrm{Hom}}\nolimits}
\def\Ext{\mathop{\mathrm{Ext}}\nolimits}
\def\SheafHom{\mathop{\mathcal{H}\!\mathit{om}}\nolimits}
\def\SheafExt{\mathop{\mathcal{E}\!\mathit{xt}}\nolimits}
\def\Sch{\mathit{Sch}}
\def\Mor{\mathop{\mathrm{Mor}}\nolimits}
\def\Ob{\mathop{\mathrm{Ob}}\nolimits}
\def\Sh{\mathop{\mathit{Sh}}\nolimits}
\def\NL{\mathop{N\!L}\nolimits}
\def\CH{\mathop{\mathrm{CH}}\nolimits}
\def\proetale{{pro\text{-}\acute{e}tale}}
\def\etale{{\acute{e}tale}}
\def\QCoh{\mathit{QCoh}}
\def\Ker{\mathop{\mathrm{Ker}}}
\def\Im{\mathop{\mathrm{Im}}}
\def\Coker{\mathop{\mathrm{Coker}}}
\def\Coim{\mathop{\mathrm{Coim}}}

% Boxtimes
%
\DeclareMathSymbol{\boxtimes}{\mathbin}{AMSa}{"02}

%
% Macros for moduli stacks/spaces
%
\def\QCohstack{\mathcal{QC}\!\mathit{oh}}
\def\Cohstack{\mathcal{C}\!\mathit{oh}}
\def\Spacesstack{\mathcal{S}\!\mathit{paces}}
\def\Quotfunctor{\mathrm{Quot}}
\def\Hilbfunctor{\mathrm{Hilb}}
\def\Curvesstack{\mathcal{C}\!\mathit{urves}}
\def\Polarizedstack{\mathcal{P}\!\mathit{olarized}}
\def\Complexesstack{\mathcal{C}\!\mathit{omplexes}}
% \Pic is the operator that assigns to X its picard group, usage \Pic(X)
% \Picardstack_{X/B} denotes the Picard stack of X over B
% \Picardfunctor_{X/B} denotes the Picard functor of X over B
\def\Pic{\mathop{\mathrm{Pic}}\nolimits}
\def\Picardstack{\mathcal{P}\!\mathit{ic}}
\def\Picardfunctor{\mathrm{Pic}}
\def\Deformationcategory{\mathcal{D}\!\mathit{ef}}
\usepackage{cleveref}
\crefname{originallemma}{original lemma}{original lemmas}
\Crefname{originallemma}{Original lemma}{Original lemmas}
\makeatletter
\newlabel{lemma-cohomology-powers-ideal-times-F}{{08AW}{source}{}{sourcepointer-1}{}}
\newlabel{lemma-cohomology-powers-ideal-times-F@cref}{{[originallemma][0][]08AW}{[1][0][]source}{}{}{}}
\makeatother
\crefname{originalreference}{original reference}{original references}
\Crefname{originalreference}{Original reference}{Original references}
\makeatletter
\newlabel{situation-formal-functions}{{08AV}{source}{}{sourcepointer-2}{}}
\newlabel{situation-formal-functions@cref}{{[originalreference][0][]08AV}{[1][0][]source}{}{}{}}
\makeatother
\title{Bounded generation of the cohomology along powers of the ideal, with a containment for the transition maps}
\author{The Stacks Project Authors}
\date{stacks-project:spaces-cohomology.tex:local-snapshot-20261012; GNU FDL 1.2 or later}
\begin{document}
\maketitle
\noindent\textit{Editorial note.} All delivered mathematical text is the author's own, copied from the pinned Stacks Project chapter file spaces-cohomology.tex; no statement or proof step was rewritten, reordered or removed.\par\medskip
\noindent\textit{Editorial note.} Preamble provenance: the pinned chapter file opens with \textbackslash{}input\{preamble\} and carries no preamble text of its own, so the project's single Stacks-wide preamble staged at sources/stacks-spaces-cohomology/preamble.tex is excerpted verbatim as three spans, lines 8-12 (verbatim package and the project's own reference/slogan/history display environments), lines 14-18 (Xy-pic setup) and lines 154-217 (theorem declarations and mathematical macros). The wrapper class and the amsmath/amssymb/amsthm/graphicx base packages are supplied by the assembler. No macro was invented or redefined.\par\medskip
\noindent\textit{Editorial note.} Same-chapter pointers are located at their \textbackslash{}label lines inside the pinned chapter file; cross-chapter pointers are located in the project's own tag list staged in the snapshot (sources/stacks-chow/tags.txt), at the verbatim lines carrying their canonical tags.\par\medskip
\noindent\textit{Allowed standard prerequisites (editorial).}\begin{itemize}
\item Finite generation of the graded module of cohomology along powers of the ideal. Same chapter spaces-cohomology.tex, label lemma-cohomology-powers-ideal-times-F (invoked at proof line 4287). It supplies finitely many elements \$x\_i \textbackslash{}in H\textasciicircum{}p(X, I\textasciicircum{}\{d\_i\}\textbackslash{}mathcal\{F\})\$ generating \$\textbackslash{}bigoplus\_\{n \textbackslash{}geq 0\} H\textasciicircum{}p(X, I\textasciicircum{}n\textbackslash{}mathcal\{F\})\$ over \$B = \textbackslash{}bigoplus\_\{n \textbackslash{}geq 0\} I\textasciicircum{}n\$, which is what produces the constant \$c = \textbackslash{}max\textbackslash{}\{d\_i\textbackslash{}\}\$.
\item The situation of formal functions: the algebraic space, the ideal sheaf and the module. Same chapter spaces-cohomology.tex, Situation \textbackslash{}ref\{situation-formal-functions\} (statement line 4274). It fixes the ambient data \$X\$, \$I\$ and \$\textbackslash{}mathcal\{F\}\$ in which both conclusions are formulated, including the graded ring \$B\$ acting on the graded module.
\end{itemize}\medskip
\ifdefined\hypertarget\hypertarget{sourcepointer-1}{}\fi
\noindent\textit{Source reference.} Original lemma 08AW: sources/stacks-spaces-cohomology/pinned-original-source.tex, source lines 4257--4257. Navigation pointer for \textbackslash{}ref\{lemma-cohomology-powers-ideal-times-F\} used inside the counted span (line 4287); the argument that consumes it is retained in the span.\par\medskip
\ifdefined\hypertarget\hypertarget{sourcepointer-2}{}\fi
\noindent\textit{Source reference.} Original reference 08AV: sources/stacks-spaces-cohomology/pinned-original-source.tex, source lines 4243--4243. Navigation pointer for \textbackslash{}ref\{situation-formal-functions\} used inside the counted span (line 4274); the argument that consumes it is retained in the span.\par\medskip
\begin{lemma}
\label{lemma-cohomology-powers-ideal-application}
In Situation \ref{situation-formal-functions}.
For every $p \geq 0$ there exists an integer $c \geq 0$ such that
\begin{enumerate}
\item the multiplication map
$I^{n - c} \otimes H^p(X, I^c\mathcal{F}) \to H^p(X, I^n\mathcal{F})$
is surjective for all $n \geq c$, and
\item the image of $H^p(X, I^{n + m}\mathcal{F}) \to H^p(X, I^n\mathcal{F})$
is contained in the submodule $I^{m - c} H^p(X, I^n\mathcal{F})$
for all $n \geq 0$, $m \geq c$.
\end{enumerate}
\end{lemma}

\begin{proof}
By Lemma \ref{lemma-cohomology-powers-ideal-times-F}
we can find $d_1, \ldots, d_t \geq 0$, and
$x_i \in H^p(X, I^{d_i}\mathcal{F})$ such that
$\bigoplus_{n \geq 0} H^p(X, I^n\mathcal{F})$ is generated
by $x_1, \ldots, x_t$ over $B = \bigoplus_{n \geq 0} I^n$.
Take $c = \max\{d_i\}$. It is clear that (1) holds.
For (2) let $b = \max(0, n - c)$.
Consider the commutative diagram of $A$-modules
$$
\xymatrix{
I^{n + m - c - b} \otimes I^b \otimes
H^p(X, I^c\mathcal{F}) \ar[r] \ar[d] &
I^{n + m - c} \otimes H^p(X, I^c\mathcal{F}) \ar[r] &
H^p(X, I^{n + m}\mathcal{F}) \ar[d] \\
I^{n + m - c - b} \otimes H^p(X, I^n\mathcal{F}) \ar[rr] & &
H^p(X, I^n\mathcal{F})
}
$$
By part (1) of the lemma the composition of the horizontal arrows
is surjective if $n + m \geq c$. On the other hand, it is clear
that $n + m - c - b \geq m - c$. Hence part (2).
\end{proof}

\end{document}
